%% kapa: exam=1 round=36 year=2025 period=2 form=A subject=appraisal_law subject_ko=감정평가관계법규 seq=17 number=17
%% source: Q-Net 감정평가사 기출문제 (cst00309) artlSeq=5231525 file=raw/1008/5231525/2025년 제36회 감정평가사 1차시험 2교시 문제지 원본.pdf
%% conversion: kordoc 4.24.51 -> tools/gen1.py (auto); review=proofread-sonnet; flags=none
\begin{question}{17}
감정평가 및 감정평가사에 관한 법령상 국토교통부장관이 감정평가법인등의 설립 인가를 반드시 취소하여야 하는 사유를 모두 고른 것은?
\begin{framedtext}
ㄱ. 감정평가법인등이 업무정지처분 기간 중에 감정평가와 관련된 상담 및 자문을 한 사실이 1차례 적발된 경우 \par
ㄴ. 국토교통부령으로 정한 감정평가법인등이 준수하여야 할 원칙과 기준을 위반하여 감정평가한 사실이 2차례 적발된 경우 \par
ㄷ. 감정평가법인이 해당 법인의 소속 감정평가사 외의 사람에게 토지등의 이용 및 개발 등에 대한 조언이나 정보제공한 사실이 3차례 적발된 경우
\end{framedtext}
\begin{choices}
  \item ㄱ
  \item ㄴ
  \item ㄱ, ㄷ
  \item ㄴ, ㄷ
  \item ㄱ, ㄴ, ㄷ
\end{choices}
\end{question}
\answer{1}
